% berlin: {"id":"synthetic-paper","published":"2026-09-09"}
\documentclass[11pt]{article}
\usepackage{amsmath}
\usepackage[hidelinks,pdfusetitle]{hyperref}
\title{A Small Study of Incremental Publishing}
\author{Example Author}
\date{September 9, 2026}
\begin{document}
\maketitle
\begin{abstract}
This synthetic paper tests conversion fidelity. It is not a research claim.
One source should retain its equations, references, and bibliography in both
the printed paper and a readable web edition.
\end{abstract}

\section{A simple cost model}\label{sec:model}
Suppose a collection has \(n\) documents and \(k\) of them change.
If rendering one document costs \(c\), a full rebuild costs
\begin{equation}\label{eq:work}
W_{\mathrm{full}}=nc,\qquad W_{\mathrm{incremental}}=kc.
\end{equation}
Equation~\eqref{eq:work} intentionally omits dependency propagation.
It provides a conversion test, not a performance prediction.\footnote{Actual
builds may also spend time on discovery, validation, indexing, and asset copying.}

\section{An illustrative comparison}\label{sec:comparison}
The model in Section~\ref{sec:model} gives the values in Table~\ref{tab:cost}.
\begin{table}[h]
\centering
\begin{tabular}{lrr}
Strategy & Documents rendered & Cost \\
\hline
Full & 100 & 100 \\
Incremental & 3 & 3 \\
\end{tabular}
\caption{Synthetic values with unit rendering cost.}\label{tab:cost}
\end{table}

The bibliography entry is explicitly fictional, used only to exercise
citation handling~\cite{example}.

\bibliographystyle{plain}
\bibliography{references}
\end{document}
